
The incompleteness argument of Rosser evaluates bounded numerical
comparisons inside \(\mathsf Q\).  The two lemmas below supply exactly
those evaluations for the object-language abbreviation \(\leq\) registered
in the hierarchy chapter, whose expansion \(r\leq t\defeq\exists
z\,(z+r=t)\) places the witness on the left of \(+\).  Both proofs
therefore use, besides the successor axioms, only the right-recursion
axioms \(x+0=x\) and \(x+Sy=S(x+y)\).

\begin{lemma}[Bounded-case enumeration in \(\mathsf Q\)]
\label{lem:q-bounded-cases}
For every \(n\in\N\),
\[
  \mathsf Q\derives\forall y\,\bigl(y\leq\bar n\to
  (y=\bar0\lor y=\bar1\lor\dots\lor y=\bar n)\bigr).
\]
\textit{Status: complete proof.}
\end{lemma}

\begin{proof}
Metainduction on \(n\); each instance is a \(\mathsf Q\)-derivation using
only the axioms \(Sx\neq0\), \(Sx=Sy\to x=y\),
\(x\neq0\to\exists y\,x=Sy\), \(x+0=x\), and \(x+Sy=S(x+y)\).

Base \(n=0\).  Reason in \(\mathsf Q\) from the expansion of
\(y\leq\bar0\): fix \(z\) with \(z+y=\bar0\).  By the case axiom,
\(y=0\) or \(y=Sw\) for some \(w\).  In the second case
\(S(z+w)=z+Sw=\bar0\) contradicts \(Sx\neq0\).  Hence \(y=\bar0\).

Step from \(n\) to \(n+1\).  Reason in \(\mathsf Q\) from \(z+y=S\bar n\),
noting that \(\overline{n+1}\) is \(S\bar n\).  If \(y=0\), the disjunct
\(y=\bar0\) holds.  If \(y=Sw\), then \(S(z+w)=z+Sw=S\bar n\), so
\(S\)-injectivity gives \(z+w=\bar n\), hence \(w\leq\bar n\); the
already-derived instance for \(n\) gives
\(w=\bar0\lor\dots\lor w=\bar n\), and applying \(S\) to each disjunct
gives \(y=\bar1\lor\dots\lor y=\overline{n+1}\).  In both cases
\(y=\bar0\lor\dots\lor y=\overline{n+1}\).
\end{proof}
